% Tapis de Quillen, with Quillen's offprint and preprint (cote 162-5, typed part): archive cote 162-5, site page 39 (PDF page 40).
% typescript. Draft compared: none.
% AI-transcribed by Claude Opus 5.5 (claude-opus-5-5, adaptive thinking, effort high), Message Batches API, batch t162_full (msgbatch_01UAh3zyWPd454AeDkhnmQ33), 2026-11-13; not yet checked by eye.
\scanpage{162-5/39}
% NOTE: the right edge of the scan cuts the ends of several typed lines (e.g. "classifica", "qu'o", "semisimplici"); end-of-line hyphens may be hidden there.
On doit trouver une approximation meilleure d'un cran à la classifica
tion topologique en prenant des  (2-Cat)-groupes resp. des (2-Cat)-groupes
abéliens,  ou plus généralement des " Gr-2-catégories "  resp.
"G\unsure{ra}b-2-catégories" . (NB  U\unsure{ne} Grab-catégorie défi\xout{\unsure{n}}nit une Grab-2catégorie
comme une catégorie définit une 2-catégorie; du point de vue semi-simplici-%
al, cette opération serait analogue à celle de "cosquélette\xout{\unsure{"}} tronqué";
Mais également une Grab-2-catégorie définie en G\unsure{r}ab-catégorie\xout{\unsure{,}} savoir la
catégorie des automorphismes de l'objet unité; cette \xout{dé} opération serait
analogue à celle de "tuage des groupes d'homotopie en dim \hand{$>$}\over{\unsure{1}}{1}").  Pour la
notion gro\xout{\unsure{s}}ssière ( Gr\xout{\illegible{1}} ou Grab) l'idée de ma lettre à Verdier était qu'o
doit trouver la même chose que les tronqués\inserted{au cran 2}, dans la catégorie \uline{homotopique}
des H-espaces  -  i.e. ce qu'on déduit des H-espaces par tuage des groupes
d'homotopie en dimensions \hand{$>$}2 .

\margin{$\|$} (Demander à Quillen relations exactes entre \xout{H}objets-groupes et objets
\xout{\unsure{ponctués}}x connexes de la catégorie homotopique\inserted[typed]{ponctuée}: forment-ils des catégories équivalen-%
tes par le foncteur espace des lacets\hand{?}).
\margin{\unsure{non}, on a un foncteur \unsure{pl. fidèle} (?) \unsure{mais non ess.} \unsure{surjectif} ....}

Pour les \uline{vrais} groupes\inserted[typed]{abéliens} de (2-Cat), l'approche de Quillen consiste à
les identifier \xout{\unsure{à}}d'abord à des groupes abéliens de la catégorie (Bicat),
donc à des groupes abéliens bi-semi\xout{e}simpliciaux particuliers (conditions
sur chaque ligne et chaque colonne). Appliquant la version bi-semisimplici
ale du tapis de Dold, on trouve que \xout{\illegible{1}} la catégorie de ces objets
estb équivalente à la catégorie des \xout{objets} bicomplexes \inserted[typed]{de chaines} de groupes  abélien
qui sont nuls en dehors du carré  \unsure{0}\hand{$\leq$}i,j \hand{$\leq$} 1.  \xout{Quillen dit que} La colonne
que la première ligne corresponde à un \xout{\unsure{ensemble}} groupe semisimplicial
trivial signifie alors qu'elle n'a que le terme de degré zéro, i.e. que
le bicomplexe est réduit à

\diagram{\begin{tikzcd}
C_{oo} & 0 \arrow[l] & \hand{$0$} \arrow[l] \\
C_{1o} \arrow[u] & C_{11} \arrow[l] \arrow[u] & \hand{$0$} \arrow[l]
\end{tikzcd}}
% NOTE: the bottom edge of the scan cuts just below the diagram; a mark under C_{1o} is partly hidden.

\mspage{39}
